\section{PLATFORM AND CONTACT MODEL}

\subsection{Mechanical and Simulation Model}

The six legs are numbered front to rear as right $1,3,5$ and left $2,4,6$.
Actuator IDs $1$--$6$ rotate the proximal links, IDs $7$--$12$ the middle links,
and IDs $13$--$18$ the distal arcs. The MuJoCo model has 57 bodies, one floating
base and 18 revolute joints, and a total modeled mass of $4.161\,$kg. Its
component masses encode measured aluminum frames and actuators plus estimated
15\%-infill polymer parts; therefore the total is a simulation parameter until
the reassembled robot is weighed. Table~\ref{tab:hardware} distinguishes that
model quantity from verified part data.

\begin{table}[t]
\caption{Actuation, contact and distal geometry used in simulation}
\label{tab:hardware}
\centering\footnotesize
\setlength{\tabcolsep}{5pt}
\begin{tabular}{@{}lll@{}}
\toprule
Quantity & Value & Role\\
\midrule
IDs 1--6 & MX-28AT, $77\,$g & proximal\\
IDs 7--12 & XM430-W350-T, $82\,$g & middle\\
IDs 13--18 & XM430-W210-T, $82\,$g & distal arc\\
FR07 / FR12 & $20/46\,$g & frames\\
Arc radii $r_i/r_o$ & $112.5/122.5\,$mm & inner/outer\\
Arc width / opening $\psi$ & $44\,$mm / $135^\circ$ & TPU contact\\
Sliding friction $\mu$ & $1.0$ & arc--ground\\
Contact softness & $3\,$mm, $\zeta=1$ & \texttt{solimp}/\texttt{solref}\\
Modeled mass & $4.161\,$kg & MuJoCo only\\
Physics / control & $500/50\,$Hz & integrate / target\\
\bottomrule
\end{tabular}
\end{table}

Each actuator is a voltage-input DC motor with position feedback rather than an
ideal torque source. For motor constants $K,R$, the outer PD loop and voltage
saturation are
\begin{align}
 \tau_i^c &= k_{p,i}(q_i^\star-q_i)+k_{d,i}(\dot q_i^\star-\dot q_i), \label{eq:pd}\\
 V_i &= \operatorname{sat}_{[-V_n,V_n]}
 \left(\frac{R_i}{K_i}\tau_i^c+K_i\dot q_i\right).            \label{eq:dcmotor}
\end{align}
Each group is specified by its rated triple (voltage, stall torque, no-load
speed), from which $K=V_n/\dot q_{\rm nl}$ and $R=KV_n/\tau_{\rm stall}$ follow:
$(12\,\mathrm{V},2.5\,\mathrm{N\,m},5.76\,\mathrm{rad/s})$,
$(12,4.1,4.82)$, and $(12,3.0,8.06)$. The simulated gains $(k_p,k_d)$ in SI units
are $(5.73,0.752)$, $(9.40,0.792)$, and $(6.88,0.264)$. Integral gain is zero;
supply voltage and a stall-matched saturation torque bound the response.

\subsection{Arc Contact and Stair-Edge Mechanics}

\begin{figure}[t]
\centering\includegraphics[width=\columnwidth]{arc_geometry.pdf}
\caption{Contact geometry of one distal link. (a) The open sector: outer and
inner radii $r_o,r_i$, opening angle $\psi$ and opening chord $l_o$, with the
lowest patch acting as the support point during walking. (b) The closed-wheel
pivot used as a reference line: the moment arm is $x_e(h)$ and the balancing
horizontal force diverges as $h\to r_o$, so any measured ascent above
$h_{\rm cw}$ must come from the opening or from articulation.}
\label{fig:geometry}
\end{figure}


Figure~\ref{fig:geometry} fixes the notation. For outer radius $r_o$, inner
radius $r_i$, and opening $\psi$, the opening chord is
\begin{equation}
 l_o=2r_o\sin(\psi/2)=0.226\ \mathrm{m}.                       \label{eq:chord}
\end{equation}
Contacts use a six-dimensional friction cone with sliding coefficient $\mu=1.0$
and a $3\,$mm compliant layer, so the TPU arc is modeled as soft but not
deformable. During walking, the support point is the centroid of mesh vertices
within $0.1\,$mm of the lowest arc patch. During rolling, this patch advances
with the distal angle.

Consider a single arc pivoting over a sharp riser of rise $h$ with its centre
initially one radius above the lower tread. The horizontal centre-to-edge offset
is
\begin{equation}
 x_e(h)=\sqrt{r_o^2-(r_o-h)^2}=\sqrt{2r_oh-h^2},\quad 0\le h\le 2r_o.
                                                               \label{eq:edgearm}
\end{equation}
If that arc carries load $F_i$, safety factor $\gamma\ge1$ and drive efficiency
$\eta\le1$ give a quasi-static lower bound on the pivot torque, while balancing
the same gravity moment with a horizontal force at the arc centre requires
\begin{align}
 \tau_{e,i}&\ge \gamma\,\frac{F_i\,x_e(h)}{\eta},               \label{eq:stairtorque}\\
 F_{h,i}&\ge \gamma\,\frac{F_i\,x_e(h)}{r_o-h},\qquad
 \mu_{\min}=\frac{|F_t|}{F_n}.                                  \label{eq:stairpush}
\end{align}
Equation~(\ref{eq:stairpush}) is the operative bound for a \emph{closed} wheel of
the same radius: $F_{h,i}\to\infty$ as $h\to r_o^-$, and the moment balance has
no solution for $h\ge r_o$. With a nosing radius $r_n$ and clearance $c$ the
closed-wheel riser limit is therefore
\begin{equation}
 h<h_{\rm cw}=r_o-r_n-c\le 122.5\ \mathrm{mm}.                  \label{eq:riserlimit}
\end{equation}
Both the $150\,$mm ($1.22r_o$) and the $200\,$mm ($1.63r_o$) risers of
Sec.~\ref{sec:results} lie above $h_{\rm cw}$, so any ascent of either one is
produced by a mechanism the closed wheel does not have. Two such mechanisms exist
on this platform, and the ablation is designed to tell them apart. The first is
the opening itself: a riser whose height is smaller than the opening chord,
$h<l_o=226\,$mm, can enter the $135^\circ$ gap and be engaged by the far end of
the C rather than pivoted over by the rim, so the arc acts as a hook. The second
is articulation, which raises and repositions the whole arc irrespective of edge
geometry. Equation~(\ref{eq:riserlimit}) is thus not a prediction of failure but
a \emph{reference line}: it converts each measured ascent above $122.5\,$mm into
evidence about which of the two mechanisms is responsible.
A second dimensionless descriptor compares the opening chord with one step's
diagonal, $l_o/\sqrt{h^2+d^2}$, giving $0.87$, $0.91$, and $0.56$ for the three
tested profiles. We report it because it is the standard legged-wheel comparison,
and note that it does not by itself order the observed outcomes.
